\section{Projective, Injective and Flat Modules}

\begin{definition}
    Let $R$ be a ring.
    \begin{itemize}
        \item A left $R$-module $P$ is called \emph{projective} if the functor
              \[
                  \Hom_R(P,-): R-\mathrm{Mod}\longrightarrow \mathbf{Ab}
              \]
              is exact.
        \item A left $R$-module $I$ is called \emph{injective} if the contravariant functor
              \[
                  \Hom_R(-,I): R-\mathrm{Mod}\longrightarrow \mathbf{Ab}
              \]
              is exact.
        \item A left $R$-module $F$ is called \emph{flat} if the functor
              \[
                  -\otimes_R F:\mathrm{Mod}\text{-}R\longrightarrow \mathbf{Ab}
              \]
              is exact.
    \end{itemize}
\end{definition}

\begin{theorem}
    Let ${}_R P$ be a left $R$-module. The following statements are equivalent.
    \begin{enumerate}
        \item Whenever we are given a commutative diagram
              \begin{center}
                  \begin{tikzcd}
                      & P \arrow[d, "\gamma"] \arrow[ld, dashed, "\widetilde{\gamma}"'] &   \\
                      M \arrow[r, two heads, "g"'] & N \arrow[r]                         & 0
                  \end{tikzcd}
              \end{center}
              with $g$ an epimorphism, there exists an $R$-homomorphism
              \[
                  \widetilde{\gamma}:P\to M
              \]
              such that $g\circ \widetilde{\gamma}=\gamma$.
        \item For every epimorphism $f:M\to N$ of left $R$-modules, the induced map
              \[
                  \Hom_R(P,f):\Hom_R(P,M)\longrightarrow \Hom_R(P,N)
              \]
              is surjective.
        \item For every ring $S$ and every right $S$-module structure on $P$ making ${}_R P_S$ an $R$-$S$-bimodule, the functor
              \[
                  \Hom_R(P,-): R - \mathrm{Mod}\longrightarrow S-\mathrm{Mod}\
              \]
              is exact.
        \item For every short exact sequence
              \[
                  0\longrightarrow M' \xrightarrow{\,u\,} M \xrightarrow{\,v\,} M'' \longrightarrow 0
              \]
              of left $R$-modules, the sequence
              \[
                  0\longrightarrow \Hom_R(P,M')
                  \xrightarrow{\,\Hom_R(P,u)\,}
                  \Hom_R(P,M)
                  \xrightarrow{\,\Hom_R(P,v)\,}
                  \Hom_R(P,M'')
                  \longrightarrow 0
              \]
              is exact.
    \end{enumerate}
\end{theorem}

\begin{proof}
    We first recall two standard facts.

    \begin{enumerate}
        \item For every left $R$-module $P$, the functor
              \[
                  \Hom_R(P,-): R\text{-}\mathrm{Mod}\to \mathbf{Ab}
              \]
              is left exact.

        \item If ${}_R P_S$ is an $R$-$S$-bimodule and $M$ is a left $R$-module, then $\Hom_R(P,M)$ is naturally a left $S$-module via
              \[
                  (s\cdot \varphi)(p):=\varphi(ps)
                  \qquad
                  (s\in S,\ \varphi\in \Hom_R(P,M),\ p\in P).
              \]
              Moreover, for every $R$-homomorphism $h:M\to N$, the induced map
              \[
                  \Hom_R(P,h):\Hom_R(P,M)\to \Hom_R(P,N)
              \]
              is $S$-linear.
    \end{enumerate}

    \medskip
    \noindent\textbf{(1) $\Longleftrightarrow$(2).}
    Let $f:M\to N$ be an epimorphism. For $\gamma\in \Hom_R(P,N)$, saying that there
    exists $\widetilde{\gamma}:P\to M$ such that
    \[
        f\circ \widetilde{\gamma}=\gamma
    \]
    is exactly saying that $\gamma$ lies in the image of
    \[
        \Hom_R(P,f):\Hom_R(P,M)\to \Hom_R(P,N).
    \]
    Hence (1) and (2) are equivalent.

    \medskip
    \noindent\textbf{(2) $\Longrightarrow$(4).}
    Let
    \[
        0\longrightarrow M' \xrightarrow{\,u\,} M \xrightarrow{\,v\,} M'' \longrightarrow 0
    \]
    be a short exact sequence of left $R$-modules. Since $\Hom_R(P,-)$ is left exact, the
    sequence
    \[
        0\longrightarrow \Hom_R(P,M')
        \xrightarrow{\,\Hom_R(P,u)\,}
        \Hom_R(P,M)
        \xrightarrow{\,\Hom_R(P,v)\,}
        \Hom_R(P,M'')
    \]
    is exact. Since $v$ is an epimorphism, (2) implies that $\Hom_R(P,v)$ is surjective.
    Therefore
    \[
        0\longrightarrow \Hom_R(P,M')
        \xrightarrow{\,\Hom_R(P,u)\,}
        \Hom_R(P,M)
        \xrightarrow{\,\Hom_R(P,v)\,}
        \Hom_R(P,M'')
        \longrightarrow 0
    \]
    is exact, proving (4).

    \medskip
    \noindent\textbf{(4) $\Longrightarrow$(2).}
    Let $f:M\to N$ be an epimorphism. Apply (4) to the short exact sequence
    \[
        0\longrightarrow \ker(f)\longrightarrow M\xrightarrow{\,f\,}N\longrightarrow 0.
    \]
    Then
    \[
        \Hom_R(P,f):\Hom_R(P,M)\to \Hom_R(P,N)
    \]
    is surjective. Thus (2) holds.

    \medskip
    \noindent\textbf{(3) $\Longrightarrow$(4).}
    Take $S=\mathbb{Z}$, with the canonical right $\mathbb{Z}$-module structure on $P$.
    Then ${}_R P_{\mathbb{Z}}$ is an $R$-$\mathbb{Z}$-bimodule, and $\mathbb{Z}\text{-}\mathrm{Mod}=\mathbf{Ab}$. Hence (3) gives that
    \[
        \Hom_R(P,-):R\text{-}\mathrm{Mod}\to \mathbf{Ab}
    \]
    is exact, which is exactly statement (4).

    \medskip
    \noindent\textbf{(4) $\Longrightarrow$(3).}
    Fix a ring $S$ and a right $S$-module structure on $P$ making ${}_R P_S$ an $R$-$S$-bimodule. Let
    \[
        0\longrightarrow M' \xrightarrow{\,u\,} M \xrightarrow{\,v\,} M'' \longrightarrow 0
    \]
    be a short exact sequence of left $R$-modules. By (4), the induced sequence
    \[
        0\longrightarrow \Hom_R(P,M')
        \xrightarrow{\,\Hom_R(P,u)\,}
        \Hom_R(P,M)
        \xrightarrow{\,\Hom_R(P,v)\,}
        \Hom_R(P,M'')
        \longrightarrow 0
    \]
    is exact as a sequence of abelian groups. By the remark above, each term carries a
    natural left $S$-module structure and each map is $S$-linear. Since exactness in $S\text{-}\mathrm{Mod}$ is equivalent to exactness of the underlying abelian groups,
    the above sequence is exact in $S\text{-}\mathrm{Mod}$. Therefore
    \[
        \Hom_R(P,-):R\text{-}\mathrm{Mod}\to S\text{-}\mathrm{Mod}
    \]
    is exact. Hence (3) holds.

    \medskip
    Therefore \textup{(1)}, \textup{(2)}, \textup{(3)}, and \textup{(4)} are equivalent.
\end{proof}

\begin{proposition}
    Free modules are projective.
\end{proposition}

\begin{proof}
    Let $F=\bigoplus_{i\in I} Re_i$ be a free left $R$-module with basis $\{e_i\}_{i\in I}$.
    Let $g:M\to N$ be an epimorphism and let $\gamma:F\to N$ be an $R$-homomorphism.
    For each $i\in I$, choose $m_i\in M$ such that
    \[
        g(m_i)=\gamma(e_i).
    \]
    Since $F$ is free, there exists a unique $R$-homomorphism
    \[
        \widetilde{\gamma}:F\to M
    \]
    such that $\widetilde{\gamma}(e_i)=m_i$ for all $i\in I$. Then
    \[
        g\widetilde{\gamma}(e_i)=g(m_i)=\gamma(e_i)
    \]
    for every basis element, so $g\widetilde{\gamma}=\gamma$. Hence $F$ is projective.
\end{proof}

\begin{theorem}
    Let ${}_R P$ be a left $R$-module. The following statements are equivalent.
    \begin{enumerate}
        \item $P$ is projective.
        \item Every epimorphism
              \[
                  f:M\longrightarrow P\longrightarrow 0
              \]
              splits, that is, there exists an $R$-homomorphism $g:P\to M$ such that
              \[
                  f\circ g=\operatorname{id}_P.
              \]
        \item $P$ is isomorphic to a direct summand of a free module.
    \end{enumerate}
\end{theorem}

\begin{proof}
    \textbf{(1) $\Longrightarrow$(2).}
    Apply the lifting property to the diagram
    \begin{center}
        \begin{tikzcd}
            & P \arrow[d, "\operatorname{id}_P"] \arrow[ld, dashed, "g"'] &   \\
            M \arrow[r, two heads, "f"'] & P \arrow[r]                     & 0
        \end{tikzcd}
    \end{center}
    to obtain $f\circ g=\operatorname{id}_P$.

    \textbf{(2) $\Longrightarrow$(3).}
    Every module is a homomorphic image of a free module, so choose an epimorphism
    \[
        \pi:F\to P
    \]
    with $F$ free. By (2), $\pi$ splits, so there exists $s:P\to F$ such that
    \[
        \pi\circ s=\operatorname{id}_P.
    \]
    Therefore
    \[
        F\cong \ker\pi \oplus s(P)\cong \ker\pi \oplus P,
    \]
    and hence $P$ is a direct summand of the free module $F$.

    \textbf{(3) $\Longrightarrow$(1).}
    Write
    \[
        F\cong P\oplus Q
    \]
    with $F$ free. Since free modules are projective, it suffices to show that a direct summand
    of a projective module is projective. Let $f:M\to N$ be an epimorphism and
    \[
        \alpha:P\to N
    \]
    an $R$-homomorphism. Define
    \[
        \widehat{\alpha}:P\oplus Q\to N,\qquad \widehat{\alpha}(p,q)=\alpha(p).
    \]
    Because $F$ is projective, there exists
    \[
        \widehat{\beta}:P\oplus Q\to M
    \]
    such that $f\circ \widehat{\beta}=\widehat{\alpha}$. Let
    \[
        \iota:P\to P\oplus Q,\qquad p\mapsto (p,0).
    \]
    Then
    \[
        f\circ (\widehat{\beta}\circ \iota)=\widehat{\alpha}\circ \iota=\alpha.
    \]
    Hence $P$ is projective.
\end{proof}

\begin{corollary}
    A left $R$-module $P$ is finitely generated and projective if and only if for some $R$-module $P'$ and some integer $n\ge 0$ there is an $R$-isomorphism
    \[
        P\oplus P'\cong R^n.
    \]
\end{corollary}

\begin{proof}
    If $P\oplus P'\cong R^n$, then $P$ is a direct summand of a finite free module, hence
    finitely generated and projective.

    Conversely, assume that $P$ is finitely generated and projective. By the theorem above,
    there exists a free module
    \[
        F=\bigoplus_{i\in I} Re_i
    \]
    and an $R$-module $Q$ such that
    \[
        F\cong P\oplus Q.
    \]
    Let $p:F\to P$ and $s:P\to F$ be the projection and inclusion corresponding to this direct
    sum decomposition. Choose generators $x_1,\dots,x_m$ of $P$. Each $s(x_j)$ is a finite
    linear combination of the basis elements $\{e_i\}_{i\in I}$, so there exists a finite
    subset $J\subseteq I$ such that
    \[
        s(x_1),\dots,s(x_m)\in F_J:=\bigoplus_{j\in J} Re_j.
    \]
    Since the $x_j$ generate $P$, we have $s(P)\subseteq F_J$. Hence the restriction
    \[
        p|_{F_J}:F_J\to P
    \]
    is surjective, and the map $s:P\to F_J$ is a section of $p|_{F_J}$. Therefore
    \[
        F_J\cong P\oplus \ker(p|_{F_J}).
    \]
    Since $F_J\cong R^{|J|}$, the result follows.
\end{proof}

